\scanpage{87}
% One long pen stroke through the leaf crosses the index and its explanatory note.
\noindent\deleted{NOTES (\textsuperscript{*})}

\diagram{\begin{tabular}{@{}r@{\enspace}p{.61\linewidth}@{\enspace}r@{\enspace}l@{}}
\deleted{1.} & \deleted{Mes amis de Survivre et Vivre} & \deleted{11} & \deleted{\handwritten{(\unsure{5|0})}} \\
\deleted{2.} & \deleted{Aldo Andreotti, Ionel Bucur} & \deleted{14} & \deleted{\handwritten{(11)}} \\
\deleted{3.} & \deleted{Jésu\deleted{\unsure{s}}s et les douze apôtres} & \deleted{25} & \deleted{\handwritten{(\unsure{19|18})}} \\
\deleted{4.} & \deleted{L'Enfant et le \deleted{m}\handwritten{M}aître} & \deleted{26} & \deleted{\handwritten{(23)}} \\
\deleted{5.} & \deleted{La peur de jouer} & \deleted{29} & \deleted{\handwritten{(23\textsuperscript{\unsure{ii|''}})}} \\
\deleted{6.} & \deleted{Les deux frères} & \deleted{29} & \deleted{\handwritten{(23\textsuperscript{\unsure{iii|'''}})}} \\
\deleted{7.} & \deleted{Causes d'un échec} & \deleted{31} & \deleted{\handwritten{(23\textsuperscript{iv})}} \\
\deleted{8.} & \deleted{Consensus déontologique - et contrôle de l'information} & \deleted{32} & \deleted{\handwritten{(25)}} \\
\deleted{9.} & \deleted{Le ``snobisme des jeunes'', ou les défenseurs de la pureté} & \deleted{33} & \deleted{\handwritten{(27)}} \\
\deleted{10.} & \deleted{Cent fers dans le feu, ou: rien ne sert de sécher !} & \deleted{36} & \deleted{\handwritten{(32)}} \\
\deleted{11.} & \deleted{L'étreinte impuissante.} & \deleted{37} & \deleted{\handwritten{(34)}} \\
\deleted{12.} & \deleted{La visi\deleted{\unsure{t|i}}te, ou le rêve malvenu} & \deleted{45} & \deleted{\handwritten{(40)}} \\
\deleted{13.} & \deleted{Krishna\deleted{\unsure{n}}murti, ou \deleted{\unsure{une}}\inserted{\handwritten{la}} libération devenue entrave} & \deleted{45} & \deleted{\handwritten{(41)}} \\
\deleted{14.} & \deleted{L'arrachement salutaire} & \deleted{45} & \deleted{\handwritten{(42)}} \\
\deleted{15.} & \deleted{Mes orphelins} & \deleted{50} & \deleted{\handwritten{(45\deleted{\unsure{g}}\unsure{a})}} \\
 & \deleted{Refus d'un héritage, ou le prix d'une contradiction\unsure{,}} & \deleted{50} & \deleted{\handwritten{(45b)}} \\
\deleted{16.} & \deleted{On n'arrête pas le progrès !} & \deleted{50} & \deleted{\handwritten{(47)}} \\
\end{tabular}}

\deleted{(\textsuperscript{*}) Voici une liste des notes les plus longues ou les plus substantielles, par ordre d'appa\unsure{ri}tion dans le texte. Les numéros placés après le \deleted{\unsure{no|nom}} nom de chaque note (lequel, je l'espère, \deleted{\unsure{les|le}}\unsure{la|les|le} rend aisément reconnaissable) indiquent celui de la section (entre 1 et 50) où la note figure, suivi de celui de la note elle-même. Une des notes, plus longue que la longue section à laquelle elle se rapporte, a été subdivisée après coup en deux parties, auxquelles il m'a paru utile de donner des noms séparés.}
% archive sheet number: 86
